\section{A periodic hitting set and the global deduction}
\label{sec:periodic}

Fix $q\in(0,1)$ for the rest of the proof. For a measurable,
$1$-periodic set $A\subseteq\R$, write
\[
 \rho(A)=m(A\cap[0,1)).
\]
We first restrict the dilation factor to $[1,2]$, while allowing every
real center. The following is the main construction.

\begin{proposition}\label{prop:periodic}
For every $p\in(0,1)$, there is an open $1$-periodic set $H\subseteq\R$
with $\rho(H)\le6p$ such that
\begin{equation}\label{eq:periodic-hitting}
 \text{for every }x\in\R\text{ and }t\in[1,2],
 \quad x+tq^n\in H\text{ for some integer }n\ge1.
\end{equation}
\end{proposition}

Sections~\ref{sec:windows}--\ref{sec:scales} prove this proposition.
We give its global consequence now, so that the subsequent construction
has a single normalized target. The summable measure budgets allow us
to use both expanding and contracting copies of the periodic set.

\begin{proof}[Proof of Theorem~\ref{thm:main}, assuming
Proposition~\ref{prop:periodic}]
For each $k\in\Z$, put
\[
 p_k=\frac{\eta}{64}\,4^{-|k|}
\]
and choose $H_k$ from Proposition~\ref{prop:periodic} with $p=p_k$.
Define
\begin{equation}\label{eq:global-open}
 C=\bigcup_{k\in\Z}\bigl(2^kH_k\cup(-2^kH_k)\bigr),
 \qquad E_{q,\eta}=[0,1]\setminus C.
\end{equation}
The set $C$ is open and symmetric under reflection, so $E_{q,\eta}$
is compact.

Each signed dilation $\pm2^kH_k$ is periodic with period $2^k$ and
has measure $2^k\rho(H_k)$ in one period. If $k\le0$, the interval
$[0,1]$ consists of exactly $2^{-k}$ periods. If $k>0$, it lies inside
a period interval of length $2^k$. In both cases,
\[
 m\bigl((\pm2^kH_k)\cap[0,1]\bigr)
 \le \max(1,2^k)\rho(H_k).
\]
Consequently,
\begin{align}
 m(C\cap[0,1])
 &\le12\sum_{k\in\Z}p_k\max(1,2^k)\notag\\
 &=\frac{12\eta}{64}
   \left(\sum_{j\ge1}4^{-j}+1+\sum_{k\ge1}2^{-k}\right)
 =\frac{7\eta}{16}<\eta.
 \label{eq:global-measure}
\end{align}
It follows that $m(E_{q,\eta})>1-\eta$.

For $s>0$, choose $k\in\Z$ and $t\in[1,2)$ with $s=2^kt$.
Apply Equation~\eqref{eq:periodic-hitting} to $H_k$ at the center
$2^{-k}x$. For some $n\ge1$ it gives
\[
 2^{-k}x+tq^n\in H_k,
 \qquad\text{hence}\qquad x+sq^n\in2^kH_k\subseteq C.
\]
For $s<0$, apply the same conclusion to $-x$ and $-s>0$ and reflect.
Thus every nontrivial affine copy of $G_q$ meets $C$ and is not
contained in $E_{q,\eta}$.
\end{proof}

The dyadic factors in Equation~\eqref{eq:global-open} only normalize
the dilation parameter. They do not change the index $n$ or require
any algebraic relation between $2$ and $q$.
